% ============================================================
%  二、模型建立
% ============================================================
\section{模型建立}

\subsection{符号与假设}

\begin{definition}[特征向量]\label{def:feature}
记样本特征 $\vb{x}=(x_1,x_2)^\top\in\R^2$，目标 $y\in\R$，设计矩阵
$\mat{X}\in\R^{n\times p}$，参数 $\vb{\beta}\in\R^p$。
\end{definition}

\begin{assumption}[线性可加噪声]
观测满足 $y_i=\vb{x}_i^\top\vb{\beta}+\varepsilon_i$，且
$\varepsilon_i\sim\mathcal{N}(0,\sigma^2)$ 独立同分布。
\end{assumption}

\subsection{回归模型}

最小二乘目标为
\begin{equation}\label{eq:ols}
  \min_{\vb{\beta}}\ \|\mat{X}\vb{\beta}-\vb{y}\|_2^2 .
\end{equation}

\begin{theorem}[正规方程]\label{thm:main}
若 $\mat{X}^\top\mat{X}$ 可逆，则 \eqref{eq:ols} 的解为
$\hat{\vb{\beta}}=(\mat{X}^\top\mat{X})^{-1}\mat{X}^\top\vb{y}$。
\end{theorem}

\begin{lemma}[数值稳定性]\label{lem:stable}
对中心化后的特征，条件数 $\kappa(\mat{X}^\top\mat{X})$ 显著下降，有利于求解稳定\upcite{smith2024synthetic}。
\end{lemma}

整体流程见 \cref{fig:workflow}；\cref{fig:timeseries} 与 \cref{fig:classification} 共同支撑结果分析。
图号区间见 \ref{fig:workflow}--\ref{fig:classification}（多图可用逗号分隔的 \verb|\cref|，见 core 文档）。

\figplot{建模流程示意}{3.6cm}{fig:workflow}{%
  运行 generate\_figures.py 生成}{fig_workflow.pdf}{%
  节点对应数据预处理、特征构建、训练与验证。}
